% Options for packages loaded elsewhere
\PassOptionsToPackage{unicode}{hyperref}
\PassOptionsToPackage{hyphens}{url}
\documentclass[
]{article}
\usepackage{xcolor}
\usepackage[margin=1in]{geometry}
\usepackage{amsmath,amssymb}
\setcounter{secnumdepth}{-\maxdimen} % remove section numbering
\usepackage{iftex}
\ifPDFTeX
  \usepackage[T1]{fontenc}
  \usepackage[utf8]{inputenc}
  \usepackage{textcomp} % provide euro and other symbols
\else % if luatex or xetex
  \usepackage{unicode-math} % this also loads fontspec
  \defaultfontfeatures{Scale=MatchLowercase}
  \defaultfontfeatures[\rmfamily]{Ligatures=TeX,Scale=1}
\fi
\usepackage{lmodern}
\ifPDFTeX\else
  % xetex/luatex font selection
\fi
% Use upquote if available, for straight quotes in verbatim environments
\IfFileExists{upquote.sty}{\usepackage{upquote}}{}
\IfFileExists{microtype.sty}{% use microtype if available
  \usepackage[]{microtype}
  \UseMicrotypeSet[protrusion]{basicmath} % disable protrusion for tt fonts
}{}
\makeatletter
\@ifundefined{KOMAClassName}{% if non-KOMA class
  \IfFileExists{parskip.sty}{%
    \usepackage{parskip}
  }{% else
    \setlength{\parindent}{0pt}
    \setlength{\parskip}{6pt plus 2pt minus 1pt}}
}{% if KOMA class
  \KOMAoptions{parskip=half}}
\makeatother
\usepackage{color}
\usepackage{fancyvrb}
\newcommand{\VerbBar}{|}
\newcommand{\VERB}{\Verb[commandchars=\\\{\}]}
\DefineVerbatimEnvironment{Highlighting}{Verbatim}{commandchars=\\\{\}}
% Add ',fontsize=\small' for more characters per line
\usepackage{framed}
\definecolor{shadecolor}{RGB}{248,248,248}
\newenvironment{Shaded}{\begin{snugshade}}{\end{snugshade}}
\newcommand{\AlertTok}[1]{\textcolor[rgb]{0.94,0.16,0.16}{#1}}
\newcommand{\AnnotationTok}[1]{\textcolor[rgb]{0.56,0.35,0.01}{\textbf{\textit{#1}}}}
\newcommand{\AttributeTok}[1]{\textcolor[rgb]{0.13,0.29,0.53}{#1}}
\newcommand{\BaseNTok}[1]{\textcolor[rgb]{0.00,0.00,0.81}{#1}}
\newcommand{\BuiltInTok}[1]{#1}
\newcommand{\CharTok}[1]{\textcolor[rgb]{0.31,0.60,0.02}{#1}}
\newcommand{\CommentTok}[1]{\textcolor[rgb]{0.56,0.35,0.01}{\textit{#1}}}
\newcommand{\CommentVarTok}[1]{\textcolor[rgb]{0.56,0.35,0.01}{\textbf{\textit{#1}}}}
\newcommand{\ConstantTok}[1]{\textcolor[rgb]{0.56,0.35,0.01}{#1}}
\newcommand{\ControlFlowTok}[1]{\textcolor[rgb]{0.13,0.29,0.53}{\textbf{#1}}}
\newcommand{\DataTypeTok}[1]{\textcolor[rgb]{0.13,0.29,0.53}{#1}}
\newcommand{\DecValTok}[1]{\textcolor[rgb]{0.00,0.00,0.81}{#1}}
\newcommand{\DocumentationTok}[1]{\textcolor[rgb]{0.56,0.35,0.01}{\textbf{\textit{#1}}}}
\newcommand{\ErrorTok}[1]{\textcolor[rgb]{0.64,0.00,0.00}{\textbf{#1}}}
\newcommand{\ExtensionTok}[1]{#1}
\newcommand{\FloatTok}[1]{\textcolor[rgb]{0.00,0.00,0.81}{#1}}
\newcommand{\FunctionTok}[1]{\textcolor[rgb]{0.13,0.29,0.53}{\textbf{#1}}}
\newcommand{\ImportTok}[1]{#1}
\newcommand{\InformationTok}[1]{\textcolor[rgb]{0.56,0.35,0.01}{\textbf{\textit{#1}}}}
\newcommand{\KeywordTok}[1]{\textcolor[rgb]{0.13,0.29,0.53}{\textbf{#1}}}
\newcommand{\NormalTok}[1]{#1}
\newcommand{\OperatorTok}[1]{\textcolor[rgb]{0.81,0.36,0.00}{\textbf{#1}}}
\newcommand{\OtherTok}[1]{\textcolor[rgb]{0.56,0.35,0.01}{#1}}
\newcommand{\PreprocessorTok}[1]{\textcolor[rgb]{0.56,0.35,0.01}{\textit{#1}}}
\newcommand{\RegionMarkerTok}[1]{#1}
\newcommand{\SpecialCharTok}[1]{\textcolor[rgb]{0.81,0.36,0.00}{\textbf{#1}}}
\newcommand{\SpecialStringTok}[1]{\textcolor[rgb]{0.31,0.60,0.02}{#1}}
\newcommand{\StringTok}[1]{\textcolor[rgb]{0.31,0.60,0.02}{#1}}
\newcommand{\VariableTok}[1]{\textcolor[rgb]{0.00,0.00,0.00}{#1}}
\newcommand{\VerbatimStringTok}[1]{\textcolor[rgb]{0.31,0.60,0.02}{#1}}
\newcommand{\WarningTok}[1]{\textcolor[rgb]{0.56,0.35,0.01}{\textbf{\textit{#1}}}}
\usepackage{graphicx}
\makeatletter
\newsavebox\pandoc@box
\newcommand*\pandocbounded[1]{% scales image to fit in text height/width
  \sbox\pandoc@box{#1}%
  \Gscale@div\@tempa{\textheight}{\dimexpr\ht\pandoc@box+\dp\pandoc@box\relax}%
  \Gscale@div\@tempb{\linewidth}{\wd\pandoc@box}%
  \ifdim\@tempb\p@<\@tempa\p@\let\@tempa\@tempb\fi% select the smaller of both
  \ifdim\@tempa\p@<\p@\scalebox{\@tempa}{\usebox\pandoc@box}%
  \else\usebox{\pandoc@box}%
  \fi%
}
% Set default figure placement to htbp
\def\fps@figure{htbp}
\makeatother
\setlength{\emergencystretch}{3em} % prevent overfull lines
\providecommand{\tightlist}{%
  \setlength{\itemsep}{0pt}\setlength{\parskip}{0pt}}
\usepackage{bookmark}
\IfFileExists{xurl.sty}{\usepackage{xurl}}{} % add URL line breaks if available
\urlstyle{same}
% fallback for those not using the hyperref driver hyperxmp:
\makeatletter
\@ifundefined{xmpquote}{\newcommand{\xmpquote}[1]{#1}}{}
\makeatother
\hypersetup{
  pdftitle={RPractical\#2},
  pdfauthor={Kent Jhon Herida},
  hidelinks,
  pdfcreator={LaTeX via pandoc}}

\title{RPractical\#2}
\author{Kent Jhon Herida}
\date{2026-11-17}

\begin{document}
\maketitle

\#\#\#Part A: Importing and inspecting a CSV.

\#A1 Import employees.csv into a data frame named emp\_df using
read.csv() with stringAsFactors = FALSE.

\begin{Shaded}
\begin{Highlighting}[]
\NormalTok{emp\_df }\OtherTok{\textless{}{-}} \FunctionTok{read.csv}\NormalTok{(}\StringTok{"employees.csv"}\NormalTok{, }\AttributeTok{stringsAsFactors =} \ConstantTok{FALSE}\NormalTok{)}
\NormalTok{emp\_df}
\end{Highlighting}
\end{Shaded}

\begin{verbatim}
##    emp_id  name department sex age salary years_service rating
## 1       1 Alice         IT   F  34  72000             8      4
## 2       2 Bruno      Sales   M  45  58000            12      3
## 3       3  Chen         IT   M  29  65000             4      5
## 4       4  Dana         HR   F  52  54000            20      3
## 5       5  Elif    Finance   F  38  78000            10      4
## 6       6 Farid      Sales   M  26  42000             2      2
## 7       7 Grace         IT   F  41  88000            15      5
## 8       8  Hiro         HR   M  33  47000             6      3
## 9       9  Ines    Finance   F  55  92000            25      4
## 10     10 Jamal      Sales   M  31  51000             5      4
## 11     11  Kara         IT   F  27  61000             3      3
## 12     12  Liam    Finance   M  48  72000            18      2
## 13     13   Mia         HR   F  36  56000             9      5
## 14     14  Noah      Sales   M  52  63000            22      4
## 15     15  Olga         IT   F  44  79000            11      4
## 16     16 Pedro    Finance   M  30  59000             5      3
## 17     17 Quinn         HR   F  25  39000             1      2
## 18     18   Raj         IT   M  50  95000            24      5
## 19     19  Sara      Sales   F  39  58000             8      3
## 20     20 Tomas    Finance   M  42  69000            14      4
\end{verbatim}

\#A2 Display the structure, the first 5 rows, and summary of emp\_df

\begin{Shaded}
\begin{Highlighting}[]
\FunctionTok{str}\NormalTok{(emp\_df)}
\end{Highlighting}
\end{Shaded}

\begin{verbatim}
## 'data.frame':    20 obs. of  8 variables:
##  $ emp_id       : int  1 2 3 4 5 6 7 8 9 10 ...
##  $ name         : chr  "Alice" "Bruno" "Chen" "Dana" ...
##  $ department   : chr  "IT" "Sales" "IT" "HR" ...
##  $ sex          : chr  "F" "M" "M" "F" ...
##  $ age          : int  34 45 29 52 38 26 41 33 55 31 ...
##  $ salary       : int  72000 58000 65000 54000 78000 42000 88000 47000 92000 51000 ...
##  $ years_service: int  8 12 4 20 10 2 15 6 25 5 ...
##  $ rating       : int  4 3 5 3 4 2 5 3 4 4 ...
\end{verbatim}

\begin{Shaded}
\begin{Highlighting}[]
\FunctionTok{head}\NormalTok{(emp\_df, }\DecValTok{5}\NormalTok{)}
\end{Highlighting}
\end{Shaded}

\begin{verbatim}
##   emp_id  name department sex age salary years_service rating
## 1      1 Alice         IT   F  34  72000             8      4
## 2      2 Bruno      Sales   M  45  58000            12      3
## 3      3  Chen         IT   M  29  65000             4      5
## 4      4  Dana         HR   F  52  54000            20      3
## 5      5  Elif    Finance   F  38  78000            10      4
\end{verbatim}

\begin{Shaded}
\begin{Highlighting}[]
\FunctionTok{summary}\NormalTok{(emp\_df)}
\end{Highlighting}
\end{Shaded}

\begin{verbatim}
##      emp_id             name        department        sex          age       
##  Min.   : 1.00   Length   :20   Length   :20   Length   :20   Min.   :25.00  
##  1st Qu.: 5.75   N.unique :20   N.unique : 4   N.unique : 2   1st Qu.:30.75  
##  Median :10.50   N.blank  : 0   N.blank  : 0   N.blank  : 0   Median :38.50  
##  Mean   :10.50   Min.nchar: 3   Min.nchar: 2   Min.nchar: 1   Mean   :38.85  
##  3rd Qu.:15.25   Max.nchar: 5   Max.nchar: 7   Max.nchar: 1   3rd Qu.:45.75  
##  Max.   :20.00                                                Max.   :55.00  
##      salary      years_service       rating   
##  Min.   :39000   Min.   : 1.00   Min.   :2.0  
##  1st Qu.:55500   1st Qu.: 5.00   1st Qu.:3.0  
##  Median :62000   Median : 9.50   Median :4.0  
##  Mean   :64900   Mean   :11.10   Mean   :3.6  
##  3rd Qu.:73500   3rd Qu.:15.75   3rd Qu.:4.0  
##  Max.   :95000   Max.   :25.00   Max.   :5.0
\end{verbatim}

\#A3 print the number of rows and the number of colunms of emp\_df,and
the class of every column

\begin{Shaded}
\begin{Highlighting}[]
\FunctionTok{nrow}\NormalTok{(emp\_df)}
\end{Highlighting}
\end{Shaded}

\begin{verbatim}
## [1] 20
\end{verbatim}

\begin{Shaded}
\begin{Highlighting}[]
\FunctionTok{ncol}\NormalTok{(emp\_df)}
\end{Highlighting}
\end{Shaded}

\begin{verbatim}
## [1] 8
\end{verbatim}

\begin{Shaded}
\begin{Highlighting}[]
\FunctionTok{sapply}\NormalTok{(emp\_df, class)}
\end{Highlighting}
\end{Shaded}

\begin{verbatim}
##        emp_id          name    department           sex           age 
##     "integer"   "character"   "character"   "character"     "integer" 
##        salary years_service        rating 
##     "integer"     "integer"     "integer"
\end{verbatim}

\#A4 Print the mean salary, the maximum age, and the number of employees
in each department using table()

\begin{Shaded}
\begin{Highlighting}[]
\FunctionTok{mean}\NormalTok{(emp\_df}\SpecialCharTok{$}\NormalTok{salary)}
\end{Highlighting}
\end{Shaded}

\begin{verbatim}
## [1] 64900
\end{verbatim}

\begin{Shaded}
\begin{Highlighting}[]
\FunctionTok{max}\NormalTok{(emp\_df}\SpecialCharTok{$}\NormalTok{age)}
\end{Highlighting}
\end{Shaded}

\begin{verbatim}
## [1] 55
\end{verbatim}

\begin{Shaded}
\begin{Highlighting}[]
\FunctionTok{table}\NormalTok{(emp\_df}\SpecialCharTok{$}\NormalTok{department)}
\end{Highlighting}
\end{Shaded}

\begin{verbatim}
## 
## Finance      HR      IT   Sales 
##       5       4       6       5
\end{verbatim}

\#\#\#B Create Data Frames

\#B1 Create a data frame new\_hires with the same column names as
emp\_df and these three employees

\begin{Shaded}
\begin{Highlighting}[]
\NormalTok{new\_hires }\OtherTok{\textless{}{-}} \FunctionTok{data.frame}\NormalTok{(}
  \AttributeTok{emp\_id =} \FunctionTok{c}\NormalTok{(}\DecValTok{21}\NormalTok{, }\DecValTok{22}\NormalTok{, }\DecValTok{23}\NormalTok{),}
  \AttributeTok{name =} \FunctionTok{c}\NormalTok{(}\StringTok{"Uma"}\NormalTok{, }\StringTok{"Victor"}\NormalTok{, }\StringTok{"Wen"}\NormalTok{),}
  \AttributeTok{department =} \FunctionTok{c}\NormalTok{(}\StringTok{"HR"}\NormalTok{, }\StringTok{"IT"}\NormalTok{, }\StringTok{"Sales"}\NormalTok{),}
  \AttributeTok{sex =} \FunctionTok{c}\NormalTok{(}\StringTok{"F"}\NormalTok{, }\StringTok{"M"}\NormalTok{, }\StringTok{"F"}\NormalTok{),}
  \AttributeTok{age =} \FunctionTok{c}\NormalTok{(}\DecValTok{28}\NormalTok{, }\DecValTok{37}\NormalTok{, }\DecValTok{33}\NormalTok{),}
  \AttributeTok{salary =} \FunctionTok{c}\NormalTok{(}\DecValTok{45000}\NormalTok{, }\DecValTok{83000}\NormalTok{, }\DecValTok{60000}\NormalTok{),}
  \AttributeTok{years\_service =} \FunctionTok{c}\NormalTok{(}\DecValTok{2}\NormalTok{, }\DecValTok{9}\NormalTok{, }\DecValTok{6}\NormalTok{),}
  \AttributeTok{rating =} \FunctionTok{c}\NormalTok{(}\DecValTok{3}\NormalTok{, }\DecValTok{4}\NormalTok{, }\DecValTok{5}\NormalTok{)}
\NormalTok{)}


\NormalTok{new\_hires}
\end{Highlighting}
\end{Shaded}

\begin{verbatim}
##   emp_id   name department sex age salary years_service rating
## 1     21    Uma         HR   F  28  45000             2      3
## 2     22 Victor         IT   M  37  83000             9      4
## 3     23    Wen      Sales   F  33  60000             6      5
\end{verbatim}

\#B2 use rbind() to append new\_hires to emp\_df, saving the result as
emp.all. print the dimensions of emp.all

\begin{Shaded}
\begin{Highlighting}[]
\NormalTok{emp\_all }\OtherTok{\textless{}{-}} \FunctionTok{rbind}\NormalTok{(emp\_df, new\_hires)}
\FunctionTok{dim}\NormalTok{(emp\_all)}
\end{Highlighting}
\end{Shaded}

\begin{verbatim}
## [1] 23  8
\end{verbatim}

\#B3 add a column bonus to emp.all equal to 5\% of salary, rounded to
the nearest whole number.

\begin{Shaded}
\begin{Highlighting}[]
\NormalTok{emp\_all}\SpecialCharTok{$}\NormalTok{bonus }\OtherTok{\textless{}{-}} \FunctionTok{round}\NormalTok{(emp\_all}\SpecialCharTok{$}\NormalTok{salary }\SpecialCharTok{*} \FloatTok{0.05}\NormalTok{)}
\FunctionTok{print}\NormalTok{(emp\_all)}
\end{Highlighting}
\end{Shaded}

\begin{verbatim}
##    emp_id   name department sex age salary years_service rating bonus
## 1       1  Alice         IT   F  34  72000             8      4  3600
## 2       2  Bruno      Sales   M  45  58000            12      3  2900
## 3       3   Chen         IT   M  29  65000             4      5  3250
## 4       4   Dana         HR   F  52  54000            20      3  2700
## 5       5   Elif    Finance   F  38  78000            10      4  3900
## 6       6  Farid      Sales   M  26  42000             2      2  2100
## 7       7  Grace         IT   F  41  88000            15      5  4400
## 8       8   Hiro         HR   M  33  47000             6      3  2350
## 9       9   Ines    Finance   F  55  92000            25      4  4600
## 10     10  Jamal      Sales   M  31  51000             5      4  2550
## 11     11   Kara         IT   F  27  61000             3      3  3050
## 12     12   Liam    Finance   M  48  72000            18      2  3600
## 13     13    Mia         HR   F  36  56000             9      5  2800
## 14     14   Noah      Sales   M  52  63000            22      4  3150
## 15     15   Olga         IT   F  44  79000            11      4  3950
## 16     16  Pedro    Finance   M  30  59000             5      3  2950
## 17     17  Quinn         HR   F  25  39000             1      2  1950
## 18     18    Raj         IT   M  50  95000            24      5  4750
## 19     19   Sara      Sales   F  39  58000             8      3  2900
## 20     20  Tomas    Finance   M  42  69000            14      4  3450
## 21     21    Uma         HR   F  28  45000             2      3  2250
## 22     22 Victor         IT   M  37  83000             9      4  4150
## 23     23    Wen      Sales   F  33  60000             6      5  3000
\end{verbatim}

\#\#\#C Subsetting

\#C1 Extract the3 first 5 rows and only the columns name and salary,
using numeric or name indexing.

\begin{Shaded}
\begin{Highlighting}[]
\NormalTok{emp\_all[}\DecValTok{1}\SpecialCharTok{:}\DecValTok{5}\NormalTok{, }\FunctionTok{c}\NormalTok{(}\StringTok{"name"}\NormalTok{, }\StringTok{"salary"}\NormalTok{)]}
\end{Highlighting}
\end{Shaded}

\begin{verbatim}
##    name salary
## 1 Alice  72000
## 2 Bruno  58000
## 3  Chen  65000
## 4  Dana  54000
## 5  Elif  78000
\end{verbatim}

\#C2 Using logical subsetting, select all employees with salary ,using
numeric or name indexing.

\begin{Shaded}
\begin{Highlighting}[]
\NormalTok{emp\_all[emp\_all}\SpecialCharTok{$}\NormalTok{salary }\SpecialCharTok{\textgreater{}} \DecValTok{70000}\NormalTok{, ]}
\end{Highlighting}
\end{Shaded}

\begin{verbatim}
##    emp_id   name department sex age salary years_service rating bonus
## 1       1  Alice         IT   F  34  72000             8      4  3600
## 5       5   Elif    Finance   F  38  78000            10      4  3900
## 7       7  Grace         IT   F  41  88000            15      5  4400
## 9       9   Ines    Finance   F  55  92000            25      4  4600
## 12     12   Liam    Finance   M  48  72000            18      2  3600
## 15     15   Olga         IT   F  44  79000            11      4  3950
## 18     18    Raj         IT   M  50  95000            24      5  4750
## 22     22 Victor         IT   M  37  83000             9      4  4150
\end{verbatim}

\#C3 using subset function, select employees in the ``IT''

\begin{Shaded}
\begin{Highlighting}[]
\FunctionTok{subset}\NormalTok{(emp\_all,}
\NormalTok{       department }\SpecialCharTok{==} \StringTok{"IT"} \SpecialCharTok{\&}\NormalTok{ years\_service }\SpecialCharTok{\textgreater{}=} \DecValTok{5}\NormalTok{,}
       \AttributeTok{select =} \FunctionTok{c}\NormalTok{(name, years\_service, salary))}
\end{Highlighting}
\end{Shaded}

\begin{verbatim}
##      name years_service salary
## 1   Alice             8  72000
## 7   Grace            15  88000
## 15   Olga            11  79000
## 18    Raj            24  95000
## 22 Victor             9  83000
\end{verbatim}

\#C4 Use which () to print the row numbers of employees older than 50,
then print their names.

\begin{Shaded}
\begin{Highlighting}[]
\FunctionTok{which}\NormalTok{(emp\_all}\SpecialCharTok{$}\NormalTok{age }\SpecialCharTok{\textgreater{}} \DecValTok{50}\NormalTok{)}
\end{Highlighting}
\end{Shaded}

\begin{verbatim}
## [1]  4  9 14
\end{verbatim}

\begin{Shaded}
\begin{Highlighting}[]
\NormalTok{emp\_all}\SpecialCharTok{$}\NormalTok{name[}\FunctionTok{which}\NormalTok{(emp\_all}\SpecialCharTok{$}\NormalTok{age }\SpecialCharTok{\textgreater{}} \DecValTok{50}\NormalTok{)]}
\end{Highlighting}
\end{Shaded}

\begin{verbatim}
## [1] "Dana" "Ines" "Noah"
\end{verbatim}

\#C5 remove the column emp\_id and bonus using negative indexing. do not
overwrite emp\_all; print the first 3 rows of the results.

\begin{Shaded}
\begin{Highlighting}[]
\NormalTok{emp\_removed }\OtherTok{\textless{}{-}}\NormalTok{ emp\_all[, }\SpecialCharTok{{-}}\FunctionTok{which}\NormalTok{(}\FunctionTok{names}\NormalTok{(emp\_all) }\SpecialCharTok{\%in\%} \FunctionTok{c}\NormalTok{(}\StringTok{"emp\_id"}\NormalTok{, }\StringTok{"bonus"}\NormalTok{))]}
\FunctionTok{head}\NormalTok{(emp\_removed, }\DecValTok{3}\NormalTok{)}
\end{Highlighting}
\end{Shaded}

\begin{verbatim}
##    name department sex age salary years_service rating
## 1 Alice         IT   F  34  72000             8      4
## 2 Bruno      Sales   M  45  58000            12      3
## 3  Chen         IT   M  29  65000             4      5
\end{verbatim}

\#C6 (2) Using logical indexing, print the salary of the employee whose
name is ``Mia''

\begin{Shaded}
\begin{Highlighting}[]
\NormalTok{emp\_all}\SpecialCharTok{$}\NormalTok{salary[emp\_all}\SpecialCharTok{$}\NormalTok{name }\SpecialCharTok{==} \StringTok{"Mia"}\NormalTok{]}
\end{Highlighting}
\end{Shaded}

\begin{verbatim}
## [1] 56000
\end{verbatim}

\#\#\#D Factors

\#D1 Convert department and sex in emp\_all to factors with factor().
Print levels() of each.

\begin{Shaded}
\begin{Highlighting}[]
\NormalTok{emp\_all}\SpecialCharTok{$}\NormalTok{department }\OtherTok{\textless{}{-}} \FunctionTok{factor}\NormalTok{(emp\_all}\SpecialCharTok{$}\NormalTok{department)}

\CommentTok{\# Convert sex into a factor}
\NormalTok{emp\_all}\SpecialCharTok{$}\NormalTok{sex }\OtherTok{\textless{}{-}} \FunctionTok{factor}\NormalTok{(emp\_all}\SpecialCharTok{$}\NormalTok{sex)}

\CommentTok{\# Print the department levels}
\FunctionTok{levels}\NormalTok{(emp\_all}\SpecialCharTok{$}\NormalTok{department)}
\end{Highlighting}
\end{Shaded}

\begin{verbatim}
## [1] "Finance" "HR"      "IT"      "Sales"
\end{verbatim}

\begin{Shaded}
\begin{Highlighting}[]
\CommentTok{\# Print the sex levels}
\FunctionTok{levels}\NormalTok{(emp\_all}\SpecialCharTok{$}\NormalTok{sex)}
\end{Highlighting}
\end{Shaded}

\begin{verbatim}
## [1] "F" "M"
\end{verbatim}

\#D2 Use table() to count employees per department, and make a two-way
table of department by 'sex.

\begin{Shaded}
\begin{Highlighting}[]
\FunctionTok{table}\NormalTok{(emp\_all}\SpecialCharTok{$}\NormalTok{department)}
\end{Highlighting}
\end{Shaded}

\begin{verbatim}
## 
## Finance      HR      IT   Sales 
##       5       5       7       6
\end{verbatim}

\#D3 Create an ordered factor rating\_level from rating using factor()
with levels = 1:5 and \#labels = c(``Poor'', ``Fair'', ``Good'', ``Very
Good'', ``Excellent''). Add it to emp\_all and print
\#table(emp\_all\$rating\_level)

\begin{Shaded}
\begin{Highlighting}[]
\NormalTok{emp\_all}\SpecialCharTok{$}\NormalTok{rating\_level }\OtherTok{\textless{}{-}} \FunctionTok{factor}\NormalTok{(}
\NormalTok{  emp\_all}\SpecialCharTok{$}\NormalTok{rating,}
  \AttributeTok{levels =} \DecValTok{1}\SpecialCharTok{:}\DecValTok{5}\NormalTok{,}
  \AttributeTok{labels =} \FunctionTok{c}\NormalTok{(}\StringTok{"Poor"}\NormalTok{, }\StringTok{"Fair"}\NormalTok{, }\StringTok{"Good"}\NormalTok{, }\StringTok{"Very Good"}\NormalTok{, }\StringTok{"Excellent"}\NormalTok{),}
  \AttributeTok{ordered =} \ConstantTok{TRUE}
\NormalTok{)}


\FunctionTok{table}\NormalTok{(emp\_all}\SpecialCharTok{$}\NormalTok{rating\_level)}
\end{Highlighting}
\end{Shaded}

\begin{verbatim}
## 
##      Poor      Fair      Good Very Good Excellent 
##         0         3         7         8         5
\end{verbatim}

\#D4 Using a for loop over levels(emp\_all\$department) and logical
subsetting, print the mean salary of \#each department, rounded to the
nearest whole number.

\begin{Shaded}
\begin{Highlighting}[]
\CommentTok{\# D4: Calculate and print the mean salary of each department using a for loop and logical subsetting}
\ControlFlowTok{for}\NormalTok{ (dept }\ControlFlowTok{in} \FunctionTok{levels}\NormalTok{(emp\_all}\SpecialCharTok{$}\NormalTok{department)) \{}
\NormalTok{  mean\_sal }\OtherTok{\textless{}{-}} \FunctionTok{round}\NormalTok{(}\FunctionTok{mean}\NormalTok{(emp\_all}\SpecialCharTok{$}\NormalTok{salary[emp\_all}\SpecialCharTok{$}\NormalTok{department }\SpecialCharTok{==}\NormalTok{ dept]))}
  \FunctionTok{print}\NormalTok{(}\FunctionTok{paste}\NormalTok{(dept, }\StringTok{":"}\NormalTok{, mean\_sal))}
\NormalTok{\}}
\end{Highlighting}
\end{Shaded}

\begin{verbatim}
## [1] "Finance : 74000"
## [1] "HR : 48200"
## [1] "IT : 77571"
## [1] "Sales : 55333"
\end{verbatim}

\#\#\#EConditional Statements

\#E1 Using nested ifelse(), add a column salary\_band to emp\_all:
``High'' if salary \textgreater= 70000, ``Medium'' \#if salary
\textgreater= 50000 (and below 70000), otherwise ``Low''. Print
table(emp\_all\$salary\_band) ..

\begin{Shaded}
\begin{Highlighting}[]
\NormalTok{emp\_all}\SpecialCharTok{$}\NormalTok{salary\_band }\OtherTok{\textless{}{-}} \FunctionTok{ifelse}\NormalTok{(}
\NormalTok{  emp\_all}\SpecialCharTok{$}\NormalTok{salary }\SpecialCharTok{\textgreater{}=} \DecValTok{70000}\NormalTok{, }\StringTok{"High"}\NormalTok{,}
  \FunctionTok{ifelse}\NormalTok{(emp\_all}\SpecialCharTok{$}\NormalTok{salary }\SpecialCharTok{\textgreater{}=} \DecValTok{50000}\NormalTok{, }\StringTok{"Medium"}\NormalTok{, }\StringTok{"Low"}\NormalTok{)}
\NormalTok{)}


\FunctionTok{table}\NormalTok{(emp\_all}\SpecialCharTok{$}\NormalTok{salary\_band)}
\end{Highlighting}
\end{Shaded}

\begin{verbatim}
## 
##   High    Low Medium 
##      8      4     11
\end{verbatim}

\#e2\# Categorize employees based on years of service

\begin{Shaded}
\begin{Highlighting}[]
\NormalTok{emp\_all}\SpecialCharTok{$}\NormalTok{level }\OtherTok{\textless{}{-}} \FunctionTok{ifelse}\NormalTok{(}
\NormalTok{  emp\_all}\SpecialCharTok{$}\NormalTok{years\_service }\SpecialCharTok{\textgreater{}=} \DecValTok{8}\NormalTok{, }\StringTok{"Senior"}\NormalTok{, }\StringTok{"Junior"}
\NormalTok{)}


\FunctionTok{table}\NormalTok{(emp\_all}\SpecialCharTok{$}\NormalTok{level)}
\end{Highlighting}
\end{Shaded}

\begin{verbatim}
## 
## Junior Senior 
##      9     14
\end{verbatim}

\#e3\# Create a function to determine the raise percentage

\begin{Shaded}
\begin{Highlighting}[]
\CommentTok{\# E3: Define the raise\_pct function}
\NormalTok{raise\_pct }\OtherTok{\textless{}{-}} \ControlFlowTok{function}\NormalTok{(rating) \{}
  \ControlFlowTok{if}\NormalTok{ (rating }\SpecialCharTok{==} \DecValTok{5}\NormalTok{) \{}
    \FunctionTok{return}\NormalTok{(}\FloatTok{0.10}\NormalTok{)}
\NormalTok{  \} }\ControlFlowTok{else} \ControlFlowTok{if}\NormalTok{ (rating }\SpecialCharTok{==} \DecValTok{4}\NormalTok{) \{}
    \FunctionTok{return}\NormalTok{(}\FloatTok{0.07}\NormalTok{)}
\NormalTok{  \} }\ControlFlowTok{else} \ControlFlowTok{if}\NormalTok{ (rating }\SpecialCharTok{==} \DecValTok{3}\NormalTok{) \{}
    \FunctionTok{return}\NormalTok{(}\FloatTok{0.04}\NormalTok{)}
\NormalTok{  \} }\ControlFlowTok{else}\NormalTok{ \{}
    \FunctionTok{return}\NormalTok{(}\DecValTok{0}\NormalTok{)}
\NormalTok{  \}}
\NormalTok{\}}


\NormalTok{emp\_all}\SpecialCharTok{$}\NormalTok{raise }\OtherTok{\textless{}{-}} \FunctionTok{sapply}\NormalTok{(emp\_all}\SpecialCharTok{$}\NormalTok{rating, raise\_pct)}


\NormalTok{emp\_all}\SpecialCharTok{$}\NormalTok{new\_salary }\OtherTok{\textless{}{-}} \FunctionTok{round}\NormalTok{(emp\_all}\SpecialCharTok{$}\NormalTok{salary }\SpecialCharTok{*}\NormalTok{ (}\DecValTok{1} \SpecialCharTok{+}\NormalTok{ emp\_all}\SpecialCharTok{$}\NormalTok{raise))}


\FunctionTok{head}\NormalTok{(emp\_all[, }\FunctionTok{c}\NormalTok{(}\StringTok{"name"}\NormalTok{, }\StringTok{"rating"}\NormalTok{, }\StringTok{"raise"}\NormalTok{, }\StringTok{"new\_salary"}\NormalTok{)], }\DecValTok{5}\NormalTok{)}
\end{Highlighting}
\end{Shaded}

\begin{verbatim}
##    name rating raise new_salary
## 1 Alice      4  0.07      77040
## 2 Bruno      3  0.04      60320
## 3  Chen      5  0.10      71500
## 4  Dana      3  0.04      56160
## 5  Elif      4  0.07      83460
\end{verbatim}

\#E4 Using an if / else statement (not ifelse()), compute total
\textless- sum(emp\_all\$new\_salary) and print \#``Budget OK'' if the
total is below 1,600,000, otherwise print ``Over budget''. Also print
the total ..

\begin{Shaded}
\begin{Highlighting}[]
\NormalTok{total }\OtherTok{\textless{}{-}} \FunctionTok{sum}\NormalTok{(emp\_all}\SpecialCharTok{$}\NormalTok{new\_salary)}


\FunctionTok{print}\NormalTok{(}\FunctionTok{paste}\NormalTok{(}\StringTok{"Total:"}\NormalTok{, total))}
\end{Highlighting}
\end{Shaded}

\begin{verbatim}
## [1] "Total: 1578770"
\end{verbatim}

\begin{Shaded}
\begin{Highlighting}[]
\ControlFlowTok{if}\NormalTok{ (total }\SpecialCharTok{\textless{}} \DecValTok{1600000}\NormalTok{) \{}
  \FunctionTok{print}\NormalTok{(}\StringTok{"Budget OK"}\NormalTok{)}
\NormalTok{\} }\ControlFlowTok{else}\NormalTok{ \{}
  \FunctionTok{print}\NormalTok{(}\StringTok{"Over budget"}\NormalTok{)}
\NormalTok{\}}
\end{Highlighting}
\end{Shaded}

\begin{verbatim}
## [1] "Budget OK"
\end{verbatim}

\#\#\#Part E: Conditional Statements

\#F1 Use sort() to print the 5 highest salary values in decreasing
order, and the first 5 employee name values \#in alphabetical order.

\begin{Shaded}
\begin{Highlighting}[]
\FunctionTok{sort}\NormalTok{(emp\_all}\SpecialCharTok{$}\NormalTok{salary, }\AttributeTok{decreasing =} \ConstantTok{TRUE}\NormalTok{)[}\DecValTok{1}\SpecialCharTok{:}\DecValTok{5}\NormalTok{]}
\end{Highlighting}
\end{Shaded}

\begin{verbatim}
## [1] 95000 92000 88000 83000 79000
\end{verbatim}

\begin{Shaded}
\begin{Highlighting}[]
\FunctionTok{sort}\NormalTok{(emp\_all}\SpecialCharTok{$}\NormalTok{name)[}\DecValTok{1}\SpecialCharTok{:}\DecValTok{5}\NormalTok{]}
\end{Highlighting}
\end{Shaded}

\begin{verbatim}
## [1] "Alice" "Bruno" "Chen"  "Dana"  "Elif"
\end{verbatim}

\#F2 Use order() to sort emp\_all by salary in descending order and
print the name, department and \#salary of the top 5. Save the sorted
data frame as 'by\_salary.

\begin{Shaded}
\begin{Highlighting}[]
\NormalTok{by\_salary }\OtherTok{\textless{}{-}}\NormalTok{ emp\_all[}\FunctionTok{order}\NormalTok{(emp\_all}\SpecialCharTok{$}\NormalTok{salary, }\AttributeTok{decreasing =} \ConstantTok{TRUE}\NormalTok{), ]}


\FunctionTok{head}\NormalTok{(by\_salary[, }\FunctionTok{c}\NormalTok{(}\StringTok{"name"}\NormalTok{, }\StringTok{"department"}\NormalTok{, }\StringTok{"salary"}\NormalTok{)], }\DecValTok{5}\NormalTok{)}
\end{Highlighting}
\end{Shaded}

\begin{verbatim}
##      name department salary
## 18    Raj         IT  95000
## 9    Ines    Finance  92000
## 7   Grace         IT  88000
## 22 Victor         IT  83000
## 15   Olga         IT  79000
\end{verbatim}

\end{document}
